
\begin{figure}[H]
\centering
\begin{tikzpicture}[
  every node/.style={figlabel},
  file/.style={statebox, minimum width=2.0cm, minimum height=.6cm, font=\sffamily\scriptsize\bfseries, align=center},
  cold/.style={neutralbox, minimum width=1.72cm, minimum height=.54cm, font=\sffamily\scriptsize, align=center}
]
\path[figpanel] (-0.55,-2.3) rectangle (11.1,3.25);
\node[figtitle, anchor=west] at (-0.1,2.82) {THE CONTEXT ENGINE};

% left context circle
\draw[BitSlate!92, line width=1.0pt] (1.55,0.25) circle (1.0cm);
\node[align=center,font=\sffamily\small] at (1.55,1.55) {Context window};
\node[align=center,figsmall] at (1.55,0.25) {ephemeral\\bounded};
\node[align=left,figsmall] at (0.86,-1.55) {-- live task context\\-- lost when the session ends};

% middle arrows
\draw[figflow, line width=1.35pt] (3.05,0.85) -- node[above, align=center, font=\sffamily\small, text=BitBlue!82!BitNavy] {PUSH OUT\\\figsmall track reality changes} (5.75,0.85);
\draw[figreturn, line width=1.35pt] (5.75,-0.35) -- node[below, align=center, font=\sffamily\small, text=BitTeal!75!black] {PULL BACK\\\figsmall load targeted state at boot} (3.05,-0.35);

% right tree
\node[anchor=west, font=\sffamily\small\bfseries] at (7.1,2.33) {State files};
\node[file] (root) at (8.55,2.0) {\statefile{STATE.md}};
\node[file] (a1) at (7.35,1.2) {\statefile{MEMORY.md}};
\node[file] (a2) at (8.55,1.2) {\statefile{TODO.md}};
\node[file] (a3) at (9.75,1.2) {\statefile{LEARNINGS.md}};
\node[cold] (c1) at (6.75,0.3) {\file{completed/}};
\node[cold] (c2) at (7.95,0.3) {\file{snapshots/}};
\node[cold] (c3) at (9.15,0.3) {\file{artifacts/}};
\node[cold] (c4) at (10.35,0.3) {\file{notes/}};
\foreach \n in {a1,a2,a3}{\draw[BitSlate!75,line width=0.9pt] (root) -- (\n);}
\foreach \a/\b in {a1/c1,a1/c2,a2/c3,a3/c4}{\draw[BitSlate!55,line width=0.8pt] (\a) -- (\b);}
\node[align=left,figsmall] at (8.55,-1.55) {-- durable\\-- selectively loaded\\-- survives across sessions};
\end{tikzpicture}
\caption{Durable state stays outside the context window. The model only pulls back the smallest useful slice.}
\label{fig:context-engine}
\end{figure}
